\documentclass[11pt,a4paper]{article}
\usepackage[margin=25mm]{geometry}
\usepackage{booktabs}
\usepackage{amsmath}
\usepackage[colorlinks=true,urlcolor=blue]{hyperref}
\title{ONYX-01: A Fictional Transforming Robot\\\large Technical Concept Dossier}
\author{Prompt Engineering Learning Project}
\date{September 2026}
\begin{document}
\maketitle
\begin{abstract}
ONYX-01 connects AI-generated visual design, mathematical modeling, control-flow diagrams, and digital presentation in a single educational prototype. All specifications are fictional educational assumptions. No physical robot or measured performance is claimed.
\end{abstract}
\section{Visual design}
The design uses layered matte obsidian armor, graphite mechanical surfaces, a pointed helmet, narrow white-blue eyes, a V-shaped chest, and shoulder wheels. A single reference image anchors subsequent images. Generated details are visual concepts rather than manufacturing drawings.
\section{Concept specifications}
\begin{center}
\begin{tabular}{ll}
\toprule
Parameter & Fictional concept value \\
\midrule
Battery energy & 2400 Wh (2.4 kWh) \\
Model speed domain & 0--60 km/h \\
Base power consumption & 200 W \\
Speed coefficient & $0.5\,\mathrm{W}/(\mathrm{km/h})^2$ \\
Low-battery threshold & 20\% \\
\bottomrule
\end{tabular}
\end{center}
\section{Energy model}
Let $v$ denote speed in km/h, $E$ energy in Wh, and $P$ power in W. Assume
\begin{align}
E &= 2400\ \mathrm{Wh},\\
P(v) &= 200 + 0.5v^2\ \mathrm{W},\\
t(v) &= \frac{E}{P(v)} = \frac{2400}{200+0.5v^2}\ \mathrm{h}, \qquad 0\leq v\leq 60.
\end{align}
The coefficient carries the units stated in the table, so both terms in $P(v)$ are powers. Energy divided by power gives time.
\begin{center}
\begin{tabular}{rrr}
\toprule
Speed (km/h) & Power (W) & Runtime (h) \\
\midrule
0 & 200 & 12 \\
20 & 400 & 6 \\
40 & 1000 & 2.4 \\
60 & 2000 & 1.2 \\
\bottomrule
\end{tabular}
\end{center}
\section{Control flow}
The conceptual controller first reads its sensors. Below 20\% battery, it stops the mission and requests charging. Otherwise, it checks for obstacles. An obstacle causes a stop and replanning before another sensor reading. With a clear path, the robot moves and repeats sensing. This is explanatory logic, not a certified controller.
\section{Prompt iteration}
The first prompt requested a short introduction. The revised prompt specified the compilation engine, document structure, two formatting examples, consistent numerical assumptions, and explicit limitations. This makes the resulting source easier to inspect and compile.
\section{Limitations}
The energy model ignores terrain, payload, temperature, changing battery efficiency, and transient motor loads. The transformation mechanism is not physically validated. AI-generated images can contain inconsistent geometry.
\end{document}
